\chapter{实时分析与语言内存模型}
\label{chap:realtime-language-memory-model}

实时正确性要求任务在截止期前完成；并发正确性要求编译器、处理器、缓存、DMA 和多个执行上下文对共享状态具有一致且被语言允许的解释。二者不能分开治理：一个存在数据竞争的“快速”任务没有可证明行为，一个使用无界锁等待的“线程安全”任务也没有可证明截止期。

本章使用经典固定优先级和 EDF 模型建立分析方法\cite{liu-layland-1973,joseph-pandya-1986}，并依据 C/C++ 抽象机和语言内存模型讨论实现约束\cite{iso-c-2024,iso-cpp-2024}。公式是设计工具而不是硬件保证；真实系统还必须测量中断、缓存、总线和设备路径。

\section{任务模型}

周期或 sporadic 任务 $\tau_i$ 常用以下参数描述：
\begin{description}
  \item[$C_i$] 最坏执行时间，包括被任务自身触发且计入其预算的代码；
  \item[$T_i$] 周期或最小到达间隔；
  \item[$D_i$] 相对截止期；
  \item[$J_i$] release jitter，即事件发生到任务进入就绪态的最大偏差；
  \item[$B_i$] 因低优先级临界区或不可抢占区域造成的最大阻塞；
  \item[$R_i$] 最坏响应时间。
\end{description}

模型必须说明时间由哪个时钟测量，任务是否自挂起，ISR、DMA completion、操作系统服务和共享中断是否计入 $C_i$。把驱动和 IRQ 时间遗漏在“系统开销”中会产生虚假的可调度结论。

\section{释放模型、执行开销与自挂起}

任务的 release 应绑定真实事件边界：定时器到期、采样完成、消息到达还是上游任务提交。事件已经发生但任务尚未被唤醒的时间属于 release jitter 或前级延迟，不能从端到端链路中消失。多个事件合并为一次唤醒时，还要说明处理的是全部积压、最新值还是有界批次。

调度器、上下文切换、tick、线程化 IRQ、cache/TLB 失效和计时插桩可以建模为任务执行时间的一部分、独立高优先级干扰或分析裕量，但必须只计一次且有测量依据。把所有开销统一填成一个未经验证的百分比，可能同时低估突发中断和高估普通路径。

自挂起会破坏许多经典分析假设。等待 DMA、总线、设备或远端核期间，任务不消耗 CPU，但醒来相位可能产生更强干扰；不能简单从 $C_i$ 中减去等待时间后套用原公式。工程上可把链路分解为 CPU segment 与有界外部等待、使用适合的自挂起分析，或通过异步状态机把等待转换为新的 sporadic release。

任务参数还可能依赖运行模式。例如冷启动、校准、故障恢复和正常稳态具有不同 $C/T/D$。应分别分析每个模式和切换瞬间，特别是旧模式工作尚未排空、新模式任务已经释放的重叠窗口。

\section{利用率与必要条件}

单核任务的总利用率为：
\begin{equation}
  U=\sum_{i=1}^{n}\frac{C_i}{T_i}
\end{equation}

$U>1$ 时单核必然不可调度；$U<1$ 只是必要条件，通常不是固定优先级系统的充分条件。对于独立、可抢占、截止期等于周期且采用 Rate Monotonic 的经典任务集，存在保守充分界：
\begin{equation}
  U \le n\left(2^{1/n}-1\right)
\end{equation}

该界随任务数增大趋近 $\ln 2$。超过该界不代表必然失败，应继续做精确响应时间分析。若存在共享资源、release jitter、非周期服务或 $D_i\ne T_i$，不能直接套用该简化判据。

\section{固定优先级响应时间分析}

对可抢占固定优先级任务，忽略更复杂自挂起时，可用迭代式计算：
\begin{equation}
  R_i^{(k+1)}=C_i+B_i+\sum_{j\in hp(i)}
  \left\lceil\frac{R_i^{(k)}+J_j}{T_j}\right\rceil C_j
\end{equation}

其中 $hp(i)$ 是优先级高于 $\tau_i$ 的任务集合。初值可以取 $C_i+B_i$，迭代到收敛；若中间结果已超过 $D_i$，该任务在当前假设下不可调度。

例如三个任务参数为 $(C,T,D)=(1,5,5)$、$(2,10,10)$ 和 $(4,20,20)$，按周期越短优先级越高。最低优先级任务的迭代为：
\begin{align}
R_3^{(0)} &= 4 \\
R_3^{(1)} &= 4+\lceil4/5\rceil\cdot1+\lceil4/10\rceil\cdot2=7 \\
R_3^{(2)} &= 4+\lceil7/5\rceil\cdot1+\lceil7/10\rceil\cdot2=8 \\
R_3^{(3)} &= 4+\lceil8/5\rceil\cdot1+\lceil8/10\rceil\cdot2=8
\end{align}
因此 $R_3=8\le D_3$。加入一个 3\,ms 的低优先级锁阻塞后，必须重新计算，不能从平均 CPU 利用率推断仍然安全。

\subsection{分析假设与敏感度}

响应时间结果应附带输入来源和敏感度。把最高优先级任务 $C_j$ 增加少量，可能因 ceiling 项多出一次抢占而让低优先级 $R_i$ 跳变；锁临界区、IRQ burst 和 release jitter 也具有类似非线性。设计评审应检查哪些参数最接近临界点，并为它们设置更严格回归门。

若迭代不收敛或超过截止期，不应只通过降低测量值“修正模型”。应考虑减少执行时间、改变优先级/周期、拆分临界区、隔离 IRQ、限制输入速率或修改架构。任何优化都要更新 WCET 与阻塞证据，平均 benchmark 改善不自动降低 $C_i$。

分析还应给出可接受裕量。$R_i$ 仅比 $D_i$ 小一个计时 tick，通常无法覆盖工具链、硬件批次和测量误差。裕量可以按绝对时间和相对比例同时报告，并由风险决定门槛，而不是把“数学上小于”视为工程完成。

\section{优先级分配}

当 $D_i=T_i$ 且满足经典假设时，Rate Monotonic 令周期越短优先级越高；当截止期不同但仍满足适用条件时，Deadline Monotonic 令相对截止期越短优先级越高。

工程系统还包括：
\begin{itemize}
  \item 硬件中断与线程化中断；
  \item 看门狗、故障检测和安全输出；
  \item 通信接收、控制计算与执行器提交；
  \item 日志、诊断、升级和后台维护。
\end{itemize}

优先级分配应从端到端截止期和故障反应时间推导。看门狗喂狗线程不应盲目设为最高优先级，否则即使控制任务死锁，它仍可能持续证明“系统健康”。

\section{优先级反转与资源协议}

高优先级任务等待低优先级任务持有的锁时发生优先级反转；若中优先级任务持续抢占锁持有者，阻塞可显著放大。

优先级继承把锁持有者临时提升到最高等待者优先级，能够限制部分反转，但不自动解决嵌套锁顺序、死锁和锁内阻塞 I/O。优先级上限类协议进一步限制任务进入临界区的条件，使阻塞项更易分析。

每个共享资源应记录：使用任务、最大临界区、嵌套关系、是否在 ISR 使用、是否可能 page fault 或访问慢设备。响应时间分析中的 $B_i$ 必须来自这些事实，而不是填零。

\section{EDF 与带宽预留}

Earliest Deadline First 动态选择绝对截止期最早的就绪任务。对独立、可抢占、单核、隐式截止期任务，$U\le1$ 是可调度的重要结论。实际系统存在调度开销、非抢占区和资源共享，仍需预留裕量并采用更完整分析。

Linux \texttt{SCHED\_DEADLINE} 使用 runtime、deadline 和 period 描述任务带宽\cite{linux-sched-deadline}。runtime 不是“期望平均时间”，而应覆盖任务在一个周期内允许消耗的 CPU；过小会被节流，过大会挤占其他服务。准入通过也不证明设备 IRQ、锁和内存路径满足端到端截止期。

非周期工作应使用有界队列、sporadic server 或带宽服务器限制影响。让诊断请求直接继承最高控制优先级，会把外部输入转换成实时拒绝服务。

\section{过载与截止期违约策略}

实际系统会遇到超出模型的输入、硬件降频、故障恢复和共享资源拥塞。实时设计必须定义过载策略，而不能只证明标称任务集可调度。可选动作包括丢弃过期样本、合并重复事件、降低非关键频率、拒绝新会话、切换简化算法或进入安全状态。

每个队列应有容量、满载策略和数据 age 上限。FIFO 保留顺序但可能处理大量已过期控制量；latest-value mailbox 保证新鲜度但丢失中间事件；覆盖计数器适合可合并事件，却不适合必须逐笔执行的事务。队列语义必须与业务不变量一致。

截止期违约应产生稳定证据：任务/链路 ID、release、start、finish、预算消耗、阻塞原因、输入 generation 和当时频率/温度。仅增加一个全局 miss counter 无法判断是执行时间增长、优先级干扰、锁等待还是时钟映射错误。

恢复策略也要防止反馈放大。任务超时后立即重试可能增加负载并让后续任务继续 miss；更合理的策略可能是取消陈旧工作、限速、释放资源并等待下一个新鲜样本。对安全功能，最大允许连续 miss 次数和进入安全状态时间应由危害分析给出。

\section{多核调度与共享干扰}

多核不能把单核利用率界简单乘以核数。partitioned 调度把任务固定到 CPU，便于在每核做响应时间分析，但可能产生装箱不足；global 调度允许迁移，利用率更灵活，却增加迁移、cache 热度和分析复杂度。项目应说明采用哪种模型以及内核实际行为。

即使任务固定到不同 CPU，仍会通过 LLC、DDR、互连、IOMMU、共享锁、设备队列和热/功耗限制互相干扰。WCET 测量必须在合法最坏共载下执行；关键核可以隔离普通任务与 IRQ，但内存控制器和系统管理中断仍可能共享。

IRQ affinity、工作队列、RCU 回调和内核线程也属于调度设计。把应用绑核而让设备 IRQ 落在同一实时 CPU，可能比不绑更差。验证应保存实际 affinity、迁移次数、频率、cache/PMU 计数和共享带宽，而不是只保存用户线程配置。

跨核锁会把远端 CPU 的执行和抢占加入本核阻塞。锁持有者若运行在可降频或高负载 CPU，优先级继承也无法消除共享硬件延迟。关键路径优先考虑静态分区、单所有者服务或消息传递，并为跨核队列定义背压和对端停止语义。

\section{端到端链路与抖动}

一个控制事务可能跨越传感器采样、DMA、ISR、任务、跨核消息、网络和执行器。只验证每个任务自身截止期，仍可能因相位和队列造成端到端超时。

链路分析应记录：
\begin{itemize}
  \item 事件时间戳所在时钟域和同步误差；
  \item 每段排队、执行、传输和设备响应上限；
  \item 数据是否使用最新值、按序消费还是允许丢弃旧样本；
  \item backpressure 时阻塞、覆盖、降采样或进入安全状态；
  \item 跨周期的 age 和 freshness，而不只记录处理耗时。
\end{itemize}

控制系统尤其应区分采样抖动、计算抖动和执行器生效抖动。平均周期稳定不能掩盖偶发双采样间隔或过期控制量。

\section{WCET 与测量边界}

WCET 可以通过静态分析、测量或混合方法获得。现代处理器的缓存、分支预测、乱序、共享总线和动态频率使严格静态界非常困难；单纯测量则只能证明“已执行样本中的最大值”。

测量型方法至少需要：
\begin{itemize}
  \item 覆盖输入、状态机路径和循环上界；
  \item 在最坏合法 CPU/DDR/外设竞争和温度下执行；
  \item 固定编译、链接、缓存、频率和内存布局；
  \item 保存最大样本对应的 trace，而非只有直方图；
  \item 为未覆盖路径、测量误差和平台变化保留裕量。
\end{itemize}

在 MCU 上，DWT cycle counter、GPIO 脉冲和 trace 可相互校验；在 Linux 上可结合 ftrace、rtla、perf 和硬件计数器。插桩自身会改变缓存和时序，关键结论应使用至少一种低侵入方法复验。

冷 cache、热 cache、首次缺页/映射、分支预测训练和稳态循环可能形成不同执行模式。只测量长循环的平均每次成本，会遗漏第一次激活或模式切换的峰值。发布门应分别覆盖冷启动、稳态、故障恢复和低功耗恢复路径。

cycle counter 需要转换到时间时，必须知道计数器频率是否恒定、是否随 DVFS 改变、是否在 suspend 停止以及是否跨核同步。GPIO 测量包含引脚写入和外部采集延迟；软件 trace 的时间戳也可能来自另一时钟域。测量报告应给出计数器语义和误差，而不只给一个微秒数字。

工具链变化会改变内联、LTO、代码布局和库调用。WCET/最大观测值应与编译器、选项、链接脚本和 binary build ID 绑定；即使源码未变，安全补丁或链接顺序变化也可能使 cache 冲突路径不同。关键实时产物升级后必须重新确认，不应无限继承旧测量。

统计分位数可以描述已观测延迟分布，但硬实时上界不能由有限样本的 P99.9 直接替代。若产品采用概率性时序目标，应预先定义样本独立性、置信度、外推方法和异常样本策略，并保留最坏样本的完整 trace。

\section{语言抽象机与未定义行为}

C/C++ 编译器只需保持语言抽象机允许观察到的行为。越界、错误生命周期、无符号/有符号混用、严格别名违规、未初始化读取、部分整数溢出和数据竞争可能产生未定义行为，编译器可以据此删除、合并或重排代码。

“汇编看起来正确”不能反向证明源程序定义良好。调试构建可用而优化构建失败时，应首先检查语言规则和同步，而不是永久关闭优化。

硬件寄存器定义还应避免实现相关位域布局。更稳健的方式是使用固定宽度整数、明确 mask/shift、静态断言和芯片厂商头文件，并限制对保留位的读改写。

\section{volatile 的边界}

\texttt{volatile} 表示每次抽象机访问都是可观察副作用，适合 MMIO、被特殊执行环境修改的对象以及少量轮询接口。它不提供：
\begin{itemize}
  \item 原子性；
  \item 多线程 happens-before；
  \item 跨核 cache 一致性；
  \item CPU 或设备总线屏障；
  \item 对复合表达式的事务性。
\end{itemize}

因此不能用 \texttt{volatile bool ready} 在线程或 hart 之间发布结构体。应使用 release store 与 acquire load，或由 RTOS/内核同步原语建立顺序。

\section{原子操作与 happens-before}

语言原子为对象提供无数据竞争访问，并通过 memory order 建立线程间顺序。典型单向发布为：
\begin{lstlisting}[language=C,caption={release/acquire 发布不可变配置}]
static struct config cfg;
static atomic_bool ready = ATOMIC_VAR_INIT(false);

void producer(void)
{
    cfg = load_and_validate_config();
    atomic_store_explicit(&ready, true, memory_order_release);
}

bool consumer(struct config *out)
{
    if (!atomic_load_explicit(&ready, memory_order_acquire))
        return false;
    *out = cfg;
    return true;
}
\end{lstlisting}

release 之前的写入在观察到该值的 acquire 之后可见。\texttt{memory\_order\_relaxed} 只保证该原子对象本身的原子性，适合独立计数等不发布其他状态的场景。

不是所有目标都能无锁实现所有原子宽度。实时路径应检查 \texttt{atomic\_is\_lock\_free()} 或工具链 ABI；库回退可能关中断、调用锁甚至需要运行库支持。

\subsection{原子生命周期与进展保证}

lock-free 只保证系统整体持续进展，不保证某个任务在有界步骤内完成；wait-free 才为每个参与者提供有界进展，但实现成本更高。基于 compare-exchange 的重试循环在竞争下可能无限失败，不应仅因“没有 mutex”就视为硬实时。

原子顺序也不能替代对象生命周期。发布指针前必须完成对象初始化，撤销后要保证所有读者退出才能释放；否则即使指针 load/store 完全原子，仍会发生 use-after-free 或 ABA。可采用静态生命周期、引用/epoch、hazard pointer、RCU 或锁，但每种回收方案都要分析内存上限与最长 grace period。

\texttt{memory\_order\_seq\_cst} 提供更强全局顺序，仍不能修复两个生产者同时写非原子 payload、设备 DMA 别名或错误 buffer 所有权。同步设计应先确定谁拥有对象和状态机，再选择满足必要顺序的最弱可审计 memory order。

原子对象的对齐和 cache line 布局影响正确性与时序。跨 cache line 或目标不支持宽度的原子可能陷入库函数；多个独立高频原子共享一条 cache line 会产生 false sharing。ABI 应静态断言 size/alignment，并在目标硬件压力下测量竞争路径。

\section{ISR、任务与多核并发}

单核 ISR 与任务并发不等同于普通线程：任务临界区可能通过屏蔽特定中断实现，ISR 则只能调用明确标记为 ISR-safe 的 API。对象的同步策略必须说明参与者是任务、ISR、另一个 hart、DMA 还是设备。

适合单生产者单消费者 ring 的 head/tail 原子，并不自动适合多生产者。无锁算法需要同时证明：
\begin{itemize}
  \item 对象生命周期和 ABA 风险；
  \item 索引回绕与容量不变量；
  \item cache line 伪共享；
  \item 目标原子指令和内存顺序；
  \item 队列满、消费者停止和复位恢复。
\end{itemize}

在安全关键路径中，一个短且可分析的锁往往优于难以验证的无锁结构。

\subsection{中断屏蔽与嵌套预算}

关中断临界区会直接增加更高优先级 IRQ 的阻塞，必须纳入 $B_i$ 或独立中断延迟预算。需要记录屏蔽的是全局中断、某个优先级阈值还是单一 IRQ；在 Cortex-M 等平台上，\texttt{PRIMASK}、\texttt{BASEPRI} 和内核临界区语义不同，不能互换估计。

嵌套中断还会消耗最坏栈空间。每层硬件自动压栈、编译器保存寄存器、FPU lazy stacking 和 RTOS 包装都要计入；只测普通调用深度无法证明异常风暴下栈安全。故障处理器本身应使用预留栈或至少验证原栈仍可信。

ISR 应只完成有界确认、时间戳和工作发布。把长协议解析、日志格式化或等待设备完成放在 ISR 中，会扩大所有低优先级任务的干扰。延后到任务上下文后，队列容量和唤醒 jitter 又成为新的分析项，必须端到端建模。

编译器屏障或 \texttt{atomic\_signal\_fence} 只约束编译器，不保证 CPU、外设或另一 hart 观察顺序。中断共享数据应使用平台支持的原子/临界区，并明确编译器与硬件两层语义。

\section{MMIO、屏障与 DMA}

语言内存模型主要描述程序对象，MMIO 和 DMA 还受处理器架构与平台影响。驱动应使用平台 accessor，例如 Linux 的 \texttt{readl()/writel()} 或 CMSIS/芯片 SDK 约定，而不是把任意地址强制转换后假定顺序成立。

屏障作用不同：
\begin{description}
  \item[编译器屏障] 限制编译器移动内存访问，不保证 CPU 或总线完成；
  \item[CPU 内存屏障] 约束处理器观察顺序，语义随架构和内存属性变化；
  \item[同步屏障] 还可能等待先前访问或 cache maintenance 完成；
  \item[设备 API] 同时封装 DMA 地址转换、一致性和所有权要求。
\end{description}

典型 DMA 状态机是 CPU owns、device owns、completion、CPU owns。所有权交接时才执行必要的 cache clean/invalidate 和屏障，禁止双方同时写同一缓冲区。描述符的 owner 位必须最后发布，设备完成标志必须先观察后读取结果字段。

平台 accessor 还可能包含端序、访问宽度和 relaxed/ordered 变体。Linux 内存屏障文档强调应优先使用内核提供的锁、原子与设备 API，而不是自行拼接架构指令\cite{linux-memory-barriers-doc}。同名 barrier 在 C、RTOS、Linux 和芯片 SDK 中未必具有相同作用域。

许多 interconnect 支持 posted write：CPU 执行写指令后，访问仍可能停留在 bridge。若软件必须确认设备已经看到配置，应按设备手册读取规定 status/flush register 或等待完成事件；随意读回同一 write-only 寄存器既可能无效，也可能有副作用。

轮询 MMIO 必须有单调时间超时、错误路径和有限 backoff。无限等待 ready bit 会把设备故障转化为系统死锁；固定次数循环则随 CPU 频率、编译优化和总线延迟变化。超时后还要决定设备是否可复位、DMA 是否仍在运行以及缓冲区所有权如何回收。

DMA completion 只说明设备或控制器达到某个协议点，不自动说明数据已经对 CPU 可见、外设物理动作已经完成或存储已持久化。驱动契约应分别定义数据可见、事务完成和物理生效边界。

\section{嵌入式 C++}

C++ 可以通过 RAII、强类型、模板和 \texttt{constexpr} 降低资源泄漏与单位错误，但项目需要明确运行时策略：
\begin{itemize}
  \item 启动前全局构造顺序和失败处理；
  \item 是否启用异常与 RTTI，以及异常表和不可恢复路径成本；
  \item 动态分配允许的阶段、内存资源和耗尽行为；
  \item 析构函数是否可能阻塞或在实时上下文执行；
  \item 标准库实现、线程 ABI 和工具链版本。
\end{itemize}

“禁用全部 C++ 特性”并非唯一策略。更有效的是按上下文建立 profile：初始化阶段可以使用较丰富设施，硬实时周期路径只允许有界、无分配且可审计的子集。

\texttt{std::atomic} 的 lock-free 属性依赖类型、对齐、架构和标准库；\texttt{std::atomic\_ref} 还要求底层对象满足生命周期与对齐条件。跨 C/C++ ABI 共享原子对象前，应验证工具链约定，而不是只比较结构体大小。

析构与 RAII 使清理路径可组合，但析构函数可能在错误展开、任务退出或锁持有期间执行。硬实时对象的析构应有界且不执行文件 I/O、等待线程或隐式释放不可预测堆；阻塞清理可通过显式 shutdown 状态机转移到非实时上下文。

\section{Rust 与 unsafe 边界}

Rust 的所有权和借用可以消除大量悬空引用与数据竞争，但 MMIO、裸指针、中断、链接脚本和 FFI 仍需要 \texttt{unsafe}。每个 unsafe 模块应写明调用者必须维护的不变量，例如地址有效、独占访问、对齐、DMA 生命周期和中断屏蔽状态。

与 C 交互时应固定 \texttt{repr(C)}、整数宽度、枚举表示、panic 策略和符号 ABI。Rust 的安全引用不能指向会被 DMA 或另一执行域在别名条件下修改的内存，除非通过适合的同步和内部可变性抽象建模。

为 MMIO、ISR 或驱动类型手工实现 \texttt{Send}/\texttt{Sync} 属于高风险 unsafe 承诺：实现者必须证明跨任务/核访问、设备副作用和对象生命周期满足要求。内部可变状态应通过 \texttt{UnsafeCell} 及受审查同步封装，不能通过裸指针绕过别名规则后仍向外暴露安全引用。

无标准库/无分配配置仍要明确 allocator 是否完全缺失、只在初始化可用还是使用固定池。panic 采用 abort、复位或自定义 handler 时，应保证关键故障证据和安全输出先完成，并验证双重 panic/栈损坏路径不会递归失控。

\section{模式切换与混合关键性}

当高关键任务超过正常预算、传感器冗余丢失或平台进入热降频时，系统可能切换到降级模式。模式切换应是一等状态机：定义触发、要停止/降频的低关键工作、预算变化、队列处理、输出限制和返回正常的证明条件。

不能在高关键任务超预算后继续让所有低关键任务按原频率运行，再期望调度器自行恢复。可以预留应急带宽、暂停诊断/界面、使用简化算法或减少数据率，但这些策略必须在最坏切换窗口分析和测试。

模式变化会改变共享数据契约。被停用生产者的最后样本是否仍有效、消费者如何识别 generation、锁由谁释放、DMA 如何取消，都要明确。安全状态需要新鲜度和有效性标志，不能把旧值静默当作当前测量。

恢复到正常模式前，应确认资源水位、温度、设备和时间同步已经稳定，并避免多个组件各自抖动切换。进入/退出阈值可使用滞回和最小驻留时间，但必须小于允许的故障暴露时间。

\section{编码规则与验证组合}

MISRA、CERT 等规则集可以减少高风险语言用法，但不能取代架构分析。推荐组合为：
\begin{enumerate}
  \item 编译器最高实用告警并将新增告警视为失败；
  \item 静态分析覆盖数据流、并发和 API 规则；
  \item 主机端 ASan/UBSan/TSan 与属性测试；
  \item 多编译器、多优化等级和 LTO 构建；
  \item 目标端边界、故障注入和实时 trace；
  \item 对偏离编码规则的情况记录理由和验证证据。
\end{enumerate}

诊断构建改变布局和时序，最终必须在发布工具链、优化和链接配置上复验功能与实时性。

\section{联合验收清单}

\begin{itemize}
  \item 每个实时任务是否具有 $C/T/D/J/B$、执行上下文和资源依赖？
  \item release、ISR、DMA、调度开销、自挂起和模式切换是否在模型中只有一个明确归属？
  \item 优先级和锁协议是否由响应时间分析而非经验数字决定？
  \item 多核 affinity、IRQ、共享 cache/DDR 和跨核锁干扰是否纳入最坏条件？
  \item 端到端链路是否验证数据 age、队列和跨时钟域误差？
  \item 过载时是否有有界丢弃、降级或安全状态，而不是无限排队和重试？
  \item 所有共享对象是否明确由语言原子、锁、关中断或消息传递保护？
  \item lock-free 算法是否同时证明对象回收、ABA、竞争下进展和最坏重试？
  \item \texttt{volatile}、屏障、MMIO/DMA API、posted write 是否分层正确，轮询与 DMA 可见/完成边界是否有界？
  \item C++ 运行时策略和 Rust unsafe/FFI 不变量是否可审查？
  \item 诊断构建是否在发布版复验，最大延迟能否关联 build ID、输入、硬件状态和 trace？
\end{itemize}

跨核内存与所有权见第~\ref{chap:heterogeneous-multicore}~章；Linux 实时与 IRQ 隔离见第~\ref{chap:realtime-linux}~章。
